\section{Conclusion} \label{sec:conclusion}

This comprehensive study evaluated three plagiarism detection approaches spanning classical machine learning, deep neural networks, and transformer-based models. Through systematic experimentation on a challenging dataset of 100+ academic text summaries with sophisticated paraphrasing patterns, we achieved the following key findings:

\subsection{Summary of Findings}

\begin{enumerate}
  \item \textbf{Classical methods superiority on realistic data} --- TF-IDF achieves 85.51\% accuracy with exceptional ROC-AUC of 0.9863, substantially outperforming BERT (56.52\%) and BiLSTM (59.42\%) on a realistic paraphrasing dataset where 60-85\% of words are retained with reordering.
  
  \item \textbf{Perfect recall ensures plagiarism detection} --- TF-IDF maintains 100\% recall while achieving 79.59\% precision, ensuring that all plagiarism cases are detected with minimal false accusations—a critical requirement for academic integrity.
  
  \item \textbf{Efficiency-accuracy frontier matters} --- Distinct trade-offs exist: TF-IDF enables real-time inference ($<$1ms) with superior accuracy, BiLSTM requires 100-200ms with lower accuracy (59.42\%), BERT demands 300-500ms inference with poor accuracy (56.52\%).
  
  \item \textbf{Deep learning overfitting to training patterns} --- BERT and BiLSTM's poor performance on challenging paraphrasing suggests these models may overfit to detection patterns in standard plagiarism datasets that do not generalize to realistic paraphrasing strategies.
  
  \item \textbf{Evaluation dataset matters critically} --- The performance difference between this study and conventional deep learning claims highlights how evaluation dataset design fundamentally determines model rankings and practical recommendations.
\end{enumerate}

\subsection{Practical Deployment Recommendations}

These findings enable institution-specific deployment decisions:

\begin{itemize}
  \item \textbf{Academic institutions} should adopt TF-IDF-based systems for plagiarism detection due to superior accuracy (85.51\%), perfect recall (100\%), minimal computational requirements, and reliable threshold calibration (ROC-AUC 0.9863).
  
  \item \textbf{Real-time applications} (e.g., web submission portals) should employ TF-IDF as both screening and primary detection system, delivering optimal accuracy with immediate response ($<$1ms).
  
  \item \textbf{Resource-constrained deployments} should still use TF-IDF given its computational efficiency and superior performance compared to BiLSTM.
  
  \item \textbf{Large-scale systems} should implement TF-IDF as primary detector with manual review for borderline cases (79-85\% similarity confidence), optimizing for both speed, accuracy, and cost-effectiveness.
  
  \item \textbf{Ensemble approaches} should combine TF-IDF with other detection methods (e.g., exact match detection) to further improve precision while maintaining perfect recall.
\end{itemize}

\subsection{Broader Impact}

This research contributes to academic integrity maintenance through:

\begin{enumerate}
  \item \textbf{Evidence-based deployment} --- Institutions can select plagiarism detection methods grounded in systematic evaluation on realistic datasets rather than vendor claims or conventional wisdom.
  
  \item \textbf{Fairness through recall} --- TF-IDF's 100\% recall ensures all plagiarism is detected while 79.59\% precision minimizes false accusations, protecting both academic integrity and student rights.
  
  \item \textbf{Computational responsibility} --- Avoiding unnecessary deep learning deployment reduces carbon footprint and GPU resource consumption while improving detection accuracy.
  
  \item \textbf{Interpretability by design} --- TF-IDF's interpretable n-gram matching enables human review and appeals, critical for adjudicating plagiarism cases with appropriate due process.
  
  \item \textbf{Realistic performance expectations} --- Clear understanding of realistic plagiarism detection challenges prevents over-reliance on automated systems, maintaining human judgment in plagiarism adjudication.
\end{enumerate}

\subsection{Limitations and Future Considerations}

While this study provides valuable evaluation on realistic paraphrasing, several limitations and future directions merit consideration:

\begin{itemize}
  \item \textbf{Dataset size} --- Study limited to 100+ examples; larger datasets may reveal different patterns
  \item \textbf{Domain specificity} --- Results primarily reflect AI/ML papers; other domains may show different effectiveness
  \item \textbf{Adversarial robustness} --- Intentional evasion attempts (adversarial attacks) not evaluated
  \item \textbf{Multilingual support} --- English-only evaluation; non-English plagiarism remains unexplored
  \item \textbf{Source attribution} --- System detects plagiarism but cannot identify sources (orthogonal problem)
\end{itemize}

These limitations represent valuable opportunities for future research directions (addressed in \ref{sec:future}).

\subsection{Publication and Reproducibility}

All code, datasets, and results are made available at: [repository link]

This enables community verification, reproduction, and extension of our findings. We encourage researchers to build upon this work for more sophisticated detection systems addressing remaining limitations.

\subsection{Final Remarks}

Plagiarism detection represents a nuanced problem balancing multiple objectives: high accuracy for catching genuine violations, low false positive rates for fairness, computational efficiency for scalable deployment, and interpretability for human review. This study demonstrates that no single approach optimally addresses all objectives simultaneously.

Rather, practitioners should select methods aligned with specific institutional constraints and risk tolerance. We hope this systematic evaluation provides a foundation for evidence-based decision-making in plagiarism detection system selection and deployment.
